
%%%%%%%%%%%%%%%%%%%%%%%%%%%%%%%%%%%%%%%%%%%%%%%%%%%%%%%%%%%
%% Inicio del resumen. 
%%%%%%%%%%%%%%%%%%%%%%%%%%%%%%%%%%%%%%%%%%%%%%%%%%%%%%%%%%%

\chapter*{Resum}

Els processos de comerç electrònic estan dissenyats per lliurar valor al client mitjançant atributs com la qualitat, el preu, els elements emocionals i socials. No obstant això, no sempre el valor modelat per l'organització correspon al valor percebut pels usuaris durant la interacció real amb el procés. Aquest desalineament pot impactar directament l'experiència del client i l'eficàcia del procés.

Aquest Treball de Fi de Màster se centra en el desenvolupament d'una eina basada en models de llenguatge (LLMs) que permeta generar i analitzar models BPMN de forma automàtica a partir de descripcions textuals o conversacionals. L'eina s'integrarà com una extensió sobre una versió local de la coneguda eina bpmn.io i permetrà, de forma iterativa, construir models de processos de negoci i identificar els valors associats al client segons la classificació PERVAL. Per avaluar l'eina desenvolupada, es dissenyarà i executarà un estudi empíric en què es compararan els models BPMN i els valors identificats automàticament amb aquells que generaria un analista expert de forma manual. L'objectiu és validar fins a quin punt l'eina és capaç d'aproximar-se a cert criteri humà tant en la construcció del model com en la identificació del valor.

Com a cas d'estudi, s'utilitzarà un procés de compra en línia inspirat en escenaris reals de comerç electrònic, sobre el qual es realitzaran diverses iteracions de modelatge i anàlisi. El principal resultat del projecte serà, per tant, doble: (1) el desenvolupament d'una eina basada en models de llenguatge per a l'anàlisi i millora iterativa de models BPMN i (2) la validació empírica i l'anàlisi de l'impacte de l'ús d'aquesta eina en la qualitat del model.

%%--------------
\newpage
%%--------------

\chapter*{Resumen}

Los procesos de comercio electrónico están diseñados para entregar valor al cliente mediante atributos como la calidad, el precio, los elementos emocionales y sociales. Sin embargo, no siempre el valor modelado por la organización corresponde al valor percibido por los usuarios durante la interacción real con el proceso. Este desalineamiento puede impactar directamente la experiencia del cliente y la eficacia del proceso.

Este Trabajo de Fin de Máster se centra en el desarrollo de una herramienta basada en modelos de lenguaje (LLMs) que permita generar y analizar modelos BPMN de forma automática a partir de descripciones textuales o conversacionales. La herramienta se integrará como una extensión sobre una versión local de la conocida herramienta bpmn.io y permitirá, de forma iterativa, construir modelos de procesos de negocio e identificar los valores asociados al cliente según la clasificación PERVAL. Para evaluar la herramienta desarrollada, se diseñará y ejecutará un estudio empírico en el que se compararán los modelos BPMN y los valores identificados automáticamente con aquellos que generaría un analista experto de forma manual. El objetivo es validar hasta qué punto la herramienta es capaz de aproximarse a cierto criterio humano tanto en la construcción del modelo como en la identificación del valor.

Como caso de estudio, se utilizará un proceso de compra en línea inspirado en escenarios reales de comercio electrónico, sobre el cual se realizarán distintas iteraciones de modelado y análisis. El principal resultado del proyecto será, por tanto, doble: (1) el desarrollo de una herramienta basada en modelos de lenguaje para el análisis y mejora iterativa de modelos BPMN y (2) la validación empírica y análisis del impacto del uso de dicha herramienta en la calidad del modelo.

%%--------------
\newpage
%%--------------


\chapter*{Abstract}

E-commerce processes are designed to deliver value to the customer through attributes such as quality, price, emotional and social elements. However, the value modelled by the organisation does not always correspond to the value perceived by users during their actual interaction with the process. This misalignment can directly impact the customer experience and the effectiveness of the process.

This Master's Thesis focuses on the development of a tool based on large language models (LLMs) that enables the automatic generation and analysis of BPMN models from textual or conversational descriptions. The tool will be integrated as an extension of a local version of the well-known bpmn.io tool and will allow, in an iterative manner, the construction of business process models and the identification of customer-associated values according to the PERVAL classification. To evaluate the developed tool, an empirical study will be designed and conducted in which the automatically generated BPMN models and identified values will be compared with those that an expert analyst would produce manually. The aim is to validate the extent to which the tool is capable of approximating human judgement both in model construction and in value identification.

As a case study, an online purchasing process inspired by real e-commerce scenarios will be used, on which several iterations of modelling and analysis will be carried out. The main outcome of the project will therefore be twofold: (1) the development of a language model-based tool for the iterative analysis and improvement of BPMN models, and (2) the empirical validation and analysis of the impact of using this tool on model quality.


%%%%%%%%%%%%%%%%%%%%%%%%%%%%%%%%%%%%%%%%%%%%%%%%%%%%%%%%%%%
%% Final del resumen. 
%%%%%%%%%%%%%%%%%%%%%%%%%%%%%%%%%%%%%%%%%%%%%%%%%%%%%%%%%%%